\answer{ex:09:1}{(a)~$P_\infty=0.2$, memória de $20\,\text{s}$. (b)~$P_\infty=1$, memória de $1\,\text{s}$. (c)~A memória é proporcional à amplitude do ruído: seis vezes.}
\answer{ex:09:2}{$P(t)=P_\infty\coth(t\sqrt{q/R})$. (a)~$\frac12\sqrt{R/q}=5\,\text{s}$, metade da memória. (b)~$t=\frac12\ln21\,\sqrt{R/q}\approx15\,\text{s}$: é quanto o \nt{ACET} elevado deve durar.}
\answer{ex:09:3}{Variância $qT/2+R/(2T)$; $T^*=\sqrt{R/q}$, variância $\sqrt{qR}=P_\infty$, como no filtro~\eqref{eq:kalman}; cresce $(\kappa+1/\kappa)/2$ vezes: $1.25$ e $5.05$.}
\answer{ex:09:4}{$a>a_w=1-\theta/(mw)$: $0.15$ com $w=0.041$, $0.22$ com $w=0.045$.}
\answer{ex:09:5}{Os erros surgem com $M\approx40$--$60$; com $M=80$, $1.9$ neurônio a mais; com $M=100$, fração do padrão $0.86$; a interferência dos outros padrões tem variância $\approx(M-1)p(1-p)/N$, de modo que $M_{\max}\approx80$, $M_{\max}\propto N/p$.}
\answer{ex:09:6}{Por exemplo, $\eta(a)=\eta\min\bigl(1,\max(0,(a-0.3)/0.6)\bigr)$. Com $\eta a$, os três neurônios alheios entraram no padrão; com $\eta(a)$, nenhum, e a escrita com $a\ge0.9$ é a mesma (fração $1.00$).}
\answer{ex:09:7}{(a)~$40$ repetições; com a $\eta(a)$ do exercício~\ref{ex:09:6}, nunca. (b)~$200$ repetições, $10$ noites.}
